\section{信道容量与信息论基础}

\subsection{信息概念的诞生}

\subsubsection{从直觉到形式化}

"明天太阳从东方升起"这句话几乎不提供信息，因为这是确定无疑的事实。"明天下雨"如果发生在久旱的沙漠，则包含了大量信息。我们的直觉告诉我们：信息量与事件的不确定性成正比。

但这一直觉需要精确化。1948年，克劳德·香农在贝尔实验室工作期间发表了一篇论文，将信息的概念数学化，奠定了信息论的基础。他的核心洞察是：信息的度量应该满足几个基本性质——信息量应为概率的减函数，确定事件的信息量为零，独立事件的联合信息量应可加。

满足这些条件的唯一形式是对数函数。对于概率为 $P(x)$ 的事件 $x$，其信息量定义为
\begin{equation}
    I(x) = -\log_2 P(x) \quad \text{比特}
\end{equation}
底数2的选择使信息量的单位是比特，即二进制位的信息量。

验证这一定义符合直觉。必然事件 $P(x)=1$ 的信息量为 $-\log_2 1 = 0$。等概率的二元事件 $P(x)=0.5$ 的信息量为 $-\log_2 0.5 = 1$ 比特。罕见事件 $P(x)=0.01$ 的信息量为 $-\log_2 0.01 \approx 6.64$ 比特，确实比常见事件携带更多信息。

\subsubsection{熵：不确定性的度量}

随机变量 $X$ 的熵定义为其平均信息量
\begin{equation}
    H(X) = -\sum_{x} P(x) \log_2 P(x)
\end{equation}
熵衡量了随机变量的不确定性。当 $X$ 服从均匀分布时，熵达到最大值；当 $X$ 确定时，熵为零。

考虑一个简单例子。设二进制信源发送1的概率为 $\alpha$，发送0的概率为 $1-\alpha$。熵为
\begin{equation}
    H(\alpha) = -\alpha \log_2 \alpha - (1-\alpha) \log_2 (1-\alpha)
\end{equation}
当 $\alpha = 0.5$ 时，$H(X) = 1$ 比特，这是最大不确定性，两种结果完全不可预测。当 $\alpha$ 偏离0.5时，熵减小；当 $\alpha = 0$ 或 $1$ 时，熵为零，结果完全确定。这一曲线呈现对称的碗形，在0.5处达到峰值。

\subsection{离散信道的容量}

\subsubsection{互信息的含义}

考虑离散信道，发送符号 $X$，接收符号 $Y$。由于噪声干扰，$Y$ 可能是 $X$ 的失真版本。我们关心的问题是：从观测到的 $Y$ 中，能获得多少关于 $X$ 的信息？

互信息回答了这个问题
\begin{equation}
    I(X;Y) = H(X) - H(X|Y)
\end{equation}
其中 $H(X|Y)$ 是条件熵，表示观察到 $Y$ 后 $X$ 的剩余不确定性。互信息衡量了观察 $Y$ 所消除的关于 $X$ 的不确定性。

互信息具有对称性，$I(X;Y) = I(Y;X) = H(Y) - H(Y|X)$。如果 $X$ 和 $Y$ 独立，$H(X|Y) = H(X)$，互信息为零，观察 $Y$ 对 $X$ 一无所知。如果 $X$ 和 $Y$ 完全相关，$H(X|Y) = 0$，互信息等于 $H(X)$，观察 $Y$ 完全确定 $X$。

\subsubsection{信道容量的定义}

信道容量是互信息的最大值，通过优化输入分布获得
\begin{equation}
    C = \max_{P(x)} I(X;Y) \quad \text{比特/符号}
\end{equation}

香农的信道编码定理指出：对于任何传输速率 $R < C$，存在编码方案使误码率任意小；对于 $R > C$，任何编码都无法实现可靠传输。容量 $C$ 是信道可靠传输的理论上限。

\subsubsection{二进制对称信道的例子}

考虑最简单的非平凡信道——二进制对称信道（BSC）。输入输出均为 $\{0,1\}$，正确传输概率为 $1-p$，错误概率为 $p$。

当输入均匀分布 $P(0)=P(1)=0.5$ 时达到容量。此时输出也均匀分布，$H(Y)=1$。条件熵 $H(Y|X) = H_b(p) = -p\log_2 p - (1-p)\log_2(1-p)$ 是二进制熵函数。因此容量为
\begin{equation}
    C = 1 - H_b(p)
\end{equation}
当 $p=0$ 时 $C=1$，信道完美传输；当 $p=0.5$ 时 $C=0$，输出与输入无关，信道完全无用。

\subsection{连续信道：香农-哈特利定理}

\subsubsection{AWGN信道模型}

加性高斯白噪声（AWGN）信道是通信理论的基本模型。接收信号是发送信号与噪声之和 $Y = X + N$，其中 $X$ 受功率约束 $E[X^2] \leq P$，$N$ 是零均值高斯噪声，方差为 $\sigma^2$。

高斯噪声的假设有其物理基础。热噪声源于导体中电子的随机热运动，根据中心极限定理，大量独立随机效应的叠加趋近高斯分布。高斯噪声还有一个数学上的便利：在相同功率约束下，高斯噪声使信道容量最小，因此以高斯噪声为假设得到的是最坏情况下的容量。

\subsubsection{香农公式的推导思路}

香农公式的完整推导需要信息论的深入工具，但其核心思想可以概述。对于带宽为 $B$ 的信号，根据采样定理，每秒需要 $2B$ 个样本。每个样本可以看作一个AWGN信道。

高斯分布使微分熵最大化。因此，使互信息最大化的输入分布也是高斯的。对于功率约束 $P$，噪声功率 $N = N_0 B$（$N_0$ 是噪声功率谱密度），每个样本的容量为 $\frac{1}{2}\log_2(1 + P/N)$。

每秒 $2B$ 个样本，总容量为
\begin{equation}
    C = B \log_2\left(1 + \frac{S}{N}\right) \quad \text{bps}
\end{equation}
这就是香农-哈特利定理，是通信工程中最重要的公式之一。

\subsubsection{香农公式的含义}

公式揭示了容量与带宽、信噪比的关系。容量与带宽 $B$ 成正比，这似乎是显然的——更宽的带宽允许更高的速率。容量与信噪比 $S/N$ 呈对数关系，这意味着信噪比每增加一倍（3 dB），容量仅增加 $B$ bps。提高信噪比的收益是递减的。

一个反直觉的结果是，带宽无限增大时容量并不无限增长。将 $N = N_0 B$ 代入，当 $B \to \infty$ 时，利用 $\ln(1+x) \approx x$ 对小 $x$ 的近似，得到
\begin{equation}
    C \to \frac{S}{N_0} \log_2 e \approx 1.44 \frac{S}{N_0}
\end{equation}
这是因为带宽增加的同时，噪声功率 $N = N_0 B$ 也线性增加，信噪比下降。存在一个由信号功率和噪声谱密度决定的基本容量极限，无法通过增加带宽来突破。

\begin{example}
电话信道带宽 $B = 3$ kHz，信噪比 $S/N = 30$ dB（即1000倍）。容量为
\begin{equation}
    C = 3000 \times \log_2(1001) \approx 30 \text{ kbps}
\end{equation}
现代调制解调器通过采用Turbo码、LDPC码等先进编码，已接近这一理论极限。
\end{example}

\subsection{衰落信道的容量}

AWGN信道是恒定的，但无线信道是时变的。信道增益 $g(t)$ 随时间随机变化，导致接收信噪比 $\gamma(t) = |g(t)|^2 P / N_0$ 也是随机的。如何定义衰落信道的容量？这取决于两个因素：信道状态信息（CSI）的可用性，以及信道变化的时标相对于编码块长度的关系。

\subsubsection{仅有接收端知道CSI}

在许多实际系统中，接收端可以通过导频信号估计信道，但发送端不知道瞬时信道状态。此时有两种重要的容量定义。

\textbf{遍历容量}适用于快速衰落场景，即信道在一个码字周期内经历多个状态。此时容量是各状态容量的统计平均
\begin{equation}
    C = \mathbb{E}\left[B \log_2(1 + \gamma)\right] = \int_0^\infty B \log_2(1 + \gamma) f_\gamma(\gamma) d\gamma
\end{equation}

对于瑞利衰落，$\gamma$ 服从指数分布 $f_\gamma(\gamma) = \frac{1}{\bar{\gamma}} e^{-\gamma/\bar{\gamma}}$。积分结果为
\begin{equation}
    \frac{C}{B} = \frac{1}{\ln 2} e^{1/\bar{\gamma}} E_1\left(\frac{1}{\bar{\gamma}}\right)
\end{equation}
其中 $E_1(x) = \int_x^\infty \frac{e^{-t}}{t} dt$ 是指数积分函数。

高信噪比近似显示，相比AWGN信道，瑞利衰落导致约2.5 dB的容量损失。这是因为衰落深谷期间的低信噪比拖累了平均性能。

\textbf{中断容量}适用于慢衰落场景，即信道在一个码字周期内近似恒定。由于发送端不知道信道状态，固定速率传输可能导致"中断"——当瞬时信道不支持该速率时，传输失败。

给定中断概率 $P_{out}$，容量定义为最大可靠传输速率
\begin{equation}
    C(P_{out}) = \max\{R: P[\log_2(1+\gamma) < R] \leq P_{out}\}
\end{equation}
对于瑞利衰落，中断概率为 $P_{out}(R) = 1 - e^{-(2^R-1)/\bar{\gamma}}$。

中断容量-中断概率曲线呈阶梯状。要求高可靠性（低 $P_{out}$）时，必须采用保守的低速率；允许较高中断概率时，可以采用更激进的高速率。这一权衡是慢衰落系统设计的核心考量。

\subsubsection{收发端都知道CSI}

当发送端也知道瞬时信道增益时，可以自适应调整发射功率和速率。这是最优策略——在深衰落时不浪费功率，在良好信道时充分利用容量。

\textbf{水填充功率分配}是这一场景下的最优策略。将功率视为"水"，信道增益的倒数视为"地面高度"，水填充算法在"地面"上倾倒固定量的水，水的高度就是分配的功率。

数学上，最优功率分配为
\begin{equation}
    P(\gamma) = \begin{cases}
        \frac{1}{\gamma_0} - \frac{1}{\gamma}, & \gamma \geq \gamma_0 \\
        0, & \gamma \lt \gamma_0
    \end{cases}
\end{equation}
其中 $\gamma_0$ 是截止信噪比，由总功率约束 $\int_{\gamma_0}^\infty P(\gamma) f_\gamma(\gamma) d\gamma = P$ 确定。

遍历容量为水填充下的平均容量
\begin{equation}
    C = B \int_{\gamma_0}^\infty \log_2\left(\frac{\gamma}{\gamma_0}\right) f_\gamma(\gamma) d\gamma
\end{equation}

水填充充分利用了信道的时间选择性。与等功率分配相比，在相同平均功率下可获得显著容量增益。这一原理的实际实现是自适应调制与编码（AMC），根据信道质量动态选择调制阶数和编码率。

\subsection{频率选择性信道的容量}

实际信道往往是频率选择性的——不同时延的多径导致不同频率经历不同衰落。这等效于一组并行的AWGN子信道
\begin{equation}
    y_k = H_k x_k + n_k, \quad k = 1, 2, ..., N
\end{equation}
其中 $H_k$ 是第 $k$ 个子信道的增益。

\textbf{频域水填充}在频域上应用同样的原理。将更多功率分配给信道增益大的频点，在深度衰落的频点减少或暂停传输
\begin{equation}
    P_k = \max\left(0, \lambda - \frac{N_0}{|H_k|^2}\right)
\end{equation}
其中 $\lambda$ 由总功率约束 $\sum_k P_k = P$ 确定。

总容量为各子信道容量之和
\begin{equation}
    C = \sum_{k=1}^N \Delta f \log_2\left(1 + \frac{P_k |H_k|^2}{N_0}\right)
\end{equation}

OFDM（正交频分复用）正是这一原理的工程实现。它将频率选择性信道转化为大量并行平坦衰落子信道，对每个子信道分别进行自适应功率和速率分配。现代WiFi、4G LTE和5G NR都基于OFDM技术。

\subsection{MIMO信道容量}

多输入多输出（MIMO）系统使用多个发射和接收天线。在丰富散射环境中，不同天线对之间的信道近似独立，MIMO信道可支持多个并行数据流。

MIMO信道模型为 $\mathbf{y} = \mathbf{H}\mathbf{x} + \mathbf{n}$，其中 $\mathbf{H}$ 是 $M_r \times M_t$ 信道矩阵。

仅接收端知道CSI时的容量为
\begin{equation}
    C = \mathbb{E}\left[\log_2 \det\left(\mathbf{I}_{M_r} + \frac{\text{SNR}}{M_t}\mathbf{H}\mathbf{H}^H\right)\right]
\end{equation}

高信噪比近似揭示了MIMO的精髓
\begin{equation}
    C \approx \min(M_t, M_r) \log_2(\text{SNR}) + \text{常数}
\end{equation}
容量随天线数的较小者线性增长。这是空间复用的理论基础——在相同带宽和功率下，通过增加天线数量成倍提高数据率。

\subsection{逼近容量的实践}

香农容量是理论极限，但工程实践不断逼近这一极限。

\textbf{信道编码}的演进是关键。卷积码在20世纪70年代实用化，距离香农限约3 dB。Turbo码于1993年提出，通过迭代译码逼近香农限至0.5 dB以内。LDPC码（最初由Gallager在1960年代提出，1990年代重新发现）具有类似的性能，且并行译码更适合硬件实现。5G NR采用LDPC作为数据信道编码，Polar码作为控制信道编码。

\textbf{自适应传输}根据信道状态动态调整调制、编码和功率。LTE和5G的AMC机制实时选择最优参数组合，使瞬时传输速率接近当前信道支持的容量。

\textbf{多天线技术}利用空间维度突破带宽限制。4G采用最多8天线，5G大规模MIMO支持64到256天线，将频谱效率提高一个数量级。

香农公式 $C = B \log_2(1 + S/N)$ 指明了通信系统设计的终极目标：在给定带宽和信噪比下，尽可能逼近这一容量界限。
